% ============================================================
%  C: Como Programar -- 6ª Edição | Deitel & Deitel
%  Capítulo 7 -- Ponteiros em C
%  FAZENDO A DIFERENÇA (7.32 e 7.33)
% ============================================================
\documentclass[12pt,a4paper]{article}
\usepackage[utf8]{inputenc}
\usepackage[T1]{fontenc}
\usepackage[brazil]{babel}
\usepackage{geometry}
\geometry{top=2.5cm,bottom=2.5cm,left=3cm,right=3cm}
\usepackage{parskip}\setlength{\parskip}{6pt}\setlength{\parindent}{0pt}
\usepackage{xcolor,tcolorbox,enumitem,listings,fancyhdr,titlesec,hyperref,booktabs,array,amsmath}
\tcbuselibrary{skins,breakable}

\definecolor{deitelblue}{RGB}{0,70,127}
\definecolor{deitelgray}{RGB}{240,242,245}
\definecolor{deitelgreen}{RGB}{0,110,60}
\definecolor{deitellightblue}{RGB}{220,235,250}
\definecolor{deitelred}{RGB}{160,20,20}
\definecolor{deitellorange}{RGB}{200,90,0}
\definecolor{ecogreen}{RGB}{0,120,50}
\definecolor{ecolightgreen}{RGB}{210,240,220}

\newtcolorbox{boaprática}[1][]{colback=deitellightblue,colframe=deitelblue,
  fonttitle=\bfseries\small,title={Boa prática de programação},breakable,#1}
\newtcolorbox{respostabox}[1][]{colback=deitelgray,colframe=deitelblue!60,
  fonttitle=\bfseries\small\color{deitelblue},title={Resposta detalhada},breakable,#1}
\newtcolorbox{enunciadoBox}[1][]{colback=white,colframe=deitelblue,
  fonttitle=\bfseries,title={#1},breakable}
\newtcolorbox{saida}[1][]{colback=black!88,colframe=deitelblue,
  fontupper=\ttfamily\small\color{green!70},
  title={\small\color{white}Saída do programa},breakable,#1}
\newtcolorbox{pesquisa}[1][]{colback=ecolightgreen,colframe=ecogreen,
  fonttitle=\bfseries\small,title={Pesquisa externa realizada},breakable,#1}
\newtcolorbox{atencao}[1][]{colback=yellow!10,colframe=deitellorange,
  fonttitle=\bfseries\small,title={Atenção},breakable,#1}

\setlist[enumerate,1]{label=\textbf{\alph*)},leftmargin=2em}
\setlist[itemize]{leftmargin=2em}

\lstset{
  basicstyle=\ttfamily\small,backgroundcolor=\color{deitelgray},
  frame=single,framerule=0.4pt,rulecolor=\color{deitelblue!40},
  breaklines=true,showstringspaces=false,
  keywordstyle=\color{deitelblue}\bfseries,
  commentstyle=\color{deitelgreen}\itshape,
  stringstyle=\color{deitelred},
  language=C,numbers=left,numberstyle=\tiny\color{gray},
  xleftmargin=18pt,xrightmargin=5pt,stepnumber=1,
}

\pagestyle{fancy}\fancyhf{}
\fancyhead[L]{\small\color{deitelblue}\textbf{C: Como Programar -- 6ª ed. | Deitel \& Deitel}}
\fancyhead[R]{\small\color{deitelblue}Cap. 7 -- Fazendo a Diferença}
\fancyfoot[C]{\small\thepage}
\renewcommand{\headrulewidth}{0.4pt}

\titleformat{\section}[block]{\large\bfseries\color{deitelblue}}{\thesection.}{0.5em}{}[\titlerule]
\titleformat{\subsection}[block]{\normalsize\bfseries\color{deitelblue!80}}{}{}{}

\hypersetup{colorlinks=true,linkcolor=deitelblue,urlcolor=ecogreen}

\begin{document}

\begin{titlepage}
  \centering\vspace*{2cm}
  {\color{deitelblue}\rule{\linewidth}{2pt}}\\[0.4cm]
  {\huge\bfseries\color{deitelblue}C: Como Programar}\\[0.2cm]
  {\large\color{deitelblue}6ª Edição $\cdot$ Paul Deitel \& Harvey Deitel}\\[0.2cm]
  {\color{deitelblue}\rule{\linewidth}{2pt}}\\[1cm]
  {\LARGE\bfseries Capítulo 7}\\[0.3cm]
  {\Large\color{deitelblue!80}Ponteiros em C}\\[0.8cm]
  \begin{tcolorbox}[colback=ecolightgreen,colframe=ecogreen,width=0.85\linewidth]
    \centering{\large\bfseries\color{ecogreen}Fazendo a Diferença}\\[0.3cm]
    {\normalsize Exercícios 7.32 e 7.33\\[0.2cm]
    Pesquisa de votação social\\
    Medição de emissão de carbono com ponteiros para função}
  \end{tcolorbox}
  \vfill{\small Abordagem incremental Deitel}
\end{titlepage}

\tableofcontents\newpage

% ============================================================
\section{Exercício 7.32 -- Pesquisa de Votação Social}
% ============================================================

\begin{enunciadoBox}[Enunciado 7.32]
Escreva um programa de votação que permita avaliar 5 questões de importância social
em escala de 1 (menos importante) a 10 (mais importante). Use um array \texttt{topics}
(\texttt{char *}) para armazenar as 5 causas e uma matriz \texttt{responses} de
5 linhas × 10 colunas para contar respostas. Apresente: relatório tabular,
média por questão, questão com maior pontuação total e questão com menor pontuação.
\end{enunciadoBox}

\subsection*{Solução em C}

\begin{respostabox}
\begin{lstlisting}
/* Ex7.32: votacao_social.c
   Sistema de pesquisa com array de ponteiros para char */
#include <stdio.h>
#define QUESTOES 5
#define NOTAS 10

int main( void )
{
    /* Array de ponteiros para strings (topicos) */
    const char *topics[ QUESTOES ] = {
        "Mudancas climaticas",
        "Educacao publica",
        "Sistema de saude",
        "Combate a pobreza",
        "Seguranca publica"
    };

    /* responses[i][j] = quantos votaram nota (j+1) para topico i */
    int responses[ QUESTOES ][ NOTAS ] = { 0 };

    int totais[ QUESTOES ] = { 0 };  /* soma de pontos por topico */
    int votantes;
    int i, j;
    int voto;
    int maiorTotal, menorTotal;
    int idxMaior, idxMenor;
    double media;

    printf( "===== PESQUISA DE OPINIAO SOCIAL =====\n" );
    printf( "Quantas pessoas vao votar? " );
    scanf( "%d", &votantes );

    /* Coleta votos */
    for ( i = 1; i <= votantes; ++i ) {
        printf( "\n--- Votante %d ---\n", i );

        for ( j = 0; j < QUESTOES; ++j ) {
            do {
                printf( "%s (1-10): ", topics[ j ] );
                scanf( "%d", &voto );
            } while ( voto < 1 || voto > 10 );

            ++responses[ j ][ voto - 1 ];  /* incrementa contador */
            totais[ j ] += voto;            /* acumula soma       */
        }
    }

    /* Cabecalho da tabela */
    printf( "\n========== RESULTADOS ==========\n" );
    printf( "%-25s", "Topico" );
    for ( j = 1; j <= NOTAS; ++j ) {
        printf( "%4d", j );
    }
    printf( "%10s%10s\n", "Total", "Media" );

    /* Item (a) e (b): relatorio tabular com totais e medias */
    for ( i = 0; i < QUESTOES; ++i ) {
        printf( "%-25s", topics[ i ] );
        for ( j = 0; j < NOTAS; ++j ) {
            printf( "%4d", responses[ i ][ j ] );
        }
        media = ( double )totais[ i ] / votantes;
        printf( "%10d%10.2f\n", totais[ i ], media );
    }

    /* Itens (c) e (d): maior e menor pontuacao */
    maiorTotal = totais[ 0 ];
    menorTotal = totais[ 0 ];
    idxMaior = 0;
    idxMenor = 0;

    for ( i = 1; i < QUESTOES; ++i ) {
        if ( totais[ i ] > maiorTotal ) {
            maiorTotal = totais[ i ];
            idxMaior = i;
        }
        if ( totais[ i ] < menorTotal ) {
            menorTotal = totais[ i ];
            idxMenor = i;
        }
    }

    printf( "\nTopico com MAIOR pontuacao total:\n" );
    printf( "  %s (total: %d pontos)\n", topics[ idxMaior ], maiorTotal );

    printf( "Topico com MENOR pontuacao total:\n" );
    printf( "  %s (total: %d pontos)\n", topics[ idxMenor ], menorTotal );

    return 0;
}
\end{lstlisting}

\textbf{Conceito-chave do Cap.~7:} usamos um \textbf{array de ponteiros para
\texttt{char}} (\texttt{const char *topics[]}). Cada elemento \texttt{topics[i]} é
um ponteiro para o primeiro caractere da string correspondente. Esta é a forma
canônica de manipular ``listas de strings'' em C -- alternativa muito mais eficiente
do que uma matriz \texttt{char topics[5][50]} (que desperdiçaria espaço).
\end{respostabox}

\subsection*{Por que isso é ``fazer a diferença''?}

\begin{tcolorbox}[colback=ecolightgreen,colframe=ecogreen,
  title={\bfseries Computação a serviço da democracia}]

Pesquisas de opinião são instrumentos fundamentais para:
\begin{itemize}[noitemsep]
  \item \textbf{Tomada de decisão democrática:} governos e ONGs identificam
  prioridades coletivas.
  \item \textbf{Participação cidadã:} pessoas expressam suas opiniões de forma
  agregada e quantificada.
  \item \textbf{Análise de tendências:} entender como prioridades mudam com o tempo.
\end{itemize}

O programa que você acabou de implementar -- embora simples -- é o núcleo
conceitual de plataformas como \textit{SurveyMonkey}, \textit{Google Forms} e
sistemas eleitorais eletrônicos. Você está construindo, em pequena escala,
ferramentas que afetam a vida cívica de milhões de pessoas.
\end{tcolorbox}

\newpage

% ============================================================
\section{Exercício 7.33 -- Emissão de Carbono com Ponteiros para Função}
% ============================================================

\begin{enunciadoBox}[Enunciado 7.33]
Crie três funções para calcular a emissão de CO\textsubscript{2} de um prédio, carro e
bicicleta. Cada função solicita dados ao usuário, calcula e exibe a emissão.
Use um array de ponteiros para função para escolher qual função executar a partir
da escolha do usuário. As funções devem retornar \texttt{void} e não receber parâmetros.
\end{enunciadoBox}

\subsection*{Pesquisa: Cálculos de Emissão de Carbono}

\begin{pesquisa}
\textbf{Fórmulas de referência (fontes: EPA, IPCC):}

\textbf{Prédios:}
\begin{itemize}[noitemsep]
  \item Consumo elétrico médio: $\approx 0{,}5$ kg CO\textsubscript{2} por kWh
  \item Consumo de gás natural: $\approx 2{,}0$ kg CO\textsubscript{2} por m³
\end{itemize}

\textbf{Carros (gasolina):}
\begin{itemize}[noitemsep]
  \item Emissão por litro de gasolina queimada: $\approx 2{,}31$ kg CO\textsubscript{2}
  \item Cálculo: $\frac{\text{km/ano}}{\text{km/litro}} \times 2{,}31$
\end{itemize}

\textbf{Bicicletas:}
\begin{itemize}[noitemsep]
  \item Emissão direta: \textbf{0}! (transporte humano, sem combustível)
  \item Emissão indireta de fabricação: $\approx 240$ kg CO\textsubscript{2} por bicicleta
  (amortizada em 5--10 anos)
  \item \textit{Compensação alimentar:} ciclistas consomem $\sim 50$ g CO\textsubscript{2}/km
  vinculadas à produção do alimento extra consumido
\end{itemize}
\end{pesquisa}

\subsection*{Solução em C}

\begin{respostabox}
\begin{lstlisting}
/* Ex7.33: emissao_carbono.c
   Array de ponteiros para funcao calculando emissoes de CO2 */
#include <stdio.h>

void emissaoPredio   ( void );
void emissaoCarro    ( void );
void emissaoBicicleta( void );

int main( void )
{
    /* Array de ponteiros para funcoes void(*)(void) */
    void ( *calculos[ 3 ] )( void ) = {
        emissaoPredio,
        emissaoCarro,
        emissaoBicicleta
    };

    int escolha;

    do {
        printf( "\n===== MEDICAO DE EMISSAO DE CO2 =====\n" );
        printf( "  0 - Predio\n" );
        printf( "  1 - Carro\n" );
        printf( "  2 - Bicicleta\n" );
        printf( "  3 - Sair\n" );
        printf( "Escolha: " );
        scanf( "%d", &escolha );

        if ( escolha >= 0 && escolha <= 2 ) {
            ( *calculos[ escolha ] )();
        }
    } while ( escolha != 3 );

    printf( "\nObrigado por se preocupar com o planeta!\n" );
    return 0;
}

void emissaoPredio( void )
{
    double consumoEletrico;  /* kWh por ano */
    double consumoGas;        /* m3 por ano */
    double emissao;

    printf( "\n--- Emissao de um predio ---\n" );
    printf( "Consumo eletrico anual (kWh): " );
    scanf( "%lf", &consumoEletrico );
    printf( "Consumo de gas natural (m3/ano): " );
    scanf( "%lf", &consumoGas );

    /* Fatores de conversao tipicos */
    emissao = consumoEletrico * 0.5 + consumoGas * 2.0;

    printf( "\n>>> Emissao anual de CO2: %.2f kg <<<\n", emissao );
    printf( "Equivalente a: %.1f arvores adultas/ano\n",
            emissao / 22.0 );  /* uma arvore absorve ~22 kg/ano */
}

void emissaoCarro( void )
{
    double kmPorAno;
    double kmPorLitro;
    double litros, emissao;

    printf( "\n--- Emissao de um carro (gasolina) ---\n" );
    printf( "Quilometros dirigidos por ano: " );
    scanf( "%lf", &kmPorAno );
    printf( "Consumo (km/litro): " );
    scanf( "%lf", &kmPorLitro );

    litros = kmPorAno / kmPorLitro;
    emissao = litros * 2.31;  /* ~2,31 kg CO2 por litro de gasolina */

    printf( "\n>>> Litros consumidos/ano: %.1f L\n", litros );
    printf( ">>> Emissao anual de CO2: %.2f kg <<<\n", emissao );
    printf( "Equivalente a: %.1f arvores adultas/ano\n",
            emissao / 22.0 );
}

void emissaoBicicleta( void )
{
    double kmPorAno;
    double emissaoFab = 240.0;  /* fabricacao, amortizada em 10 anos */
    double emissaoUso, emissaoTotal;

    printf( "\n--- Emissao de uma bicicleta ---\n" );
    printf( "Quilometros pedalados por ano: " );
    scanf( "%lf", &kmPorAno );

    /* Emissao indireta da producao de alimento extra (~50 g/km) */
    emissaoUso = kmPorAno * 0.05;  /* kg de CO2 */
    emissaoTotal = emissaoFab / 10.0 + emissaoUso;

    printf( "\n>>> Emissao direta do uso: ZERO!\n" );
    printf( ">>> Emissao indireta (fabricacao/10 anos + alimento):\n" );
    printf( "    %.2f kg de CO2/ano\n", emissaoTotal );

    /* Comparativo */
    printf( "\nComparativo: voce evitaria aproximadamente\n" );
    printf( "  %.0f kg de CO2/ano usando bicicleta em vez de carro\n",
            ( kmPorAno / 10.0 ) * 2.31 - emissaoTotal );
}
\end{lstlisting}
\end{respostabox}

\subsection*{Análise didática}

\begin{boaprática}
\textbf{Por que ponteiros para função aqui?} Note que poderíamos resolver o
problema com um simples \texttt{switch}. Mas usando ponteiros para função:

\begin{itemize}[noitemsep]
  \item Adicionar nova categoria (avião, comida, etc.) requer apenas: definir a
  função + adicioná-la ao array. Sem mexer na lógica principal!
  \item Demonstra \textbf{polimorfismo} em C: todas as funções têm o mesmo
  ``contrato'' (\texttt{void f(void)}), podendo ser chamadas de maneira uniforme.
  \item Prepara para conceitos avançados: \textit{callbacks}, \textit{event
  handlers}, \textit{strategy pattern} (orientação a objetos no Cap.~15+).
\end{itemize}
\end{boaprática}

\bigskip

\begin{atencao}
Os valores numéricos usados são aproximações didáticas. Cálculos profissionais
de pegada de carbono levam em conta dezenas de fatores adicionais: emissões
``embutidas'' em construção, manufatura, transporte de matérias-primas, etc.
Ferramentas como a \textit{WWF Footprint Calculator} ou \textit{EPA Carbon Calculator}
oferecem precisão muito maior.
\end{atencao}

\subsection*{Por que isso é ``fazer a diferença''?}

\begin{tcolorbox}[colback=ecolightgreen,colframe=ecogreen,
  title={\bfseries Conexão com a sustentabilidade global}]

Este exercício se conecta diretamente aos Exercícios 1.10 (calculadora de emissão)
e 3.47 (crescimento populacional) -- todos parte de um arco didático maior:
usar programação para entender e mitigar os desafios ambientais do século XXI.

\textbf{Ações concretas que esse programa permite:}
\begin{itemize}[noitemsep]
  \item Indivíduos avaliam sua pegada e identificam oportunidades de redução
  \item Famílias comparam transporte público vs.\ carro vs.\ bicicleta
  \item Empresas calculam emissões de operações e definem metas
  \item Educadores ilustram o conceito de \textit{decisões cotidianas com impacto
  acumulado em escala planetária}
\end{itemize}

\textbf{Programação não é uma atividade neutra:} as ferramentas que você cria
podem reforçar ou enfraquecer hábitos sustentáveis. Cada calculadora, app de
mapas ou recomendador de produtos toma micro-decisões que, agregadas, moldam
comportamentos sociais.
\end{tcolorbox}

\newpage

% ============================================================
\section*{Reflexão Final do Capítulo 7}

\begin{tcolorbox}[colback=deitellightblue,colframe=deitelblue,
  title={\bfseries Ponteiros: o poder e a responsabilidade}]

O Capítulo~7 marca um divisor de águas em sua jornada com C. Ponteiros são, ao
mesmo tempo:

\begin{itemize}[noitemsep]
  \item \textbf{O recurso mais poderoso} -- permitem manipulação direta de memória,
  passagem eficiente de grandes estruturas, polimorfismo simples e estruturas
  de dados sofisticadas (Cap.~12).
  \item \textbf{O recurso mais perigoso} -- ponteiros perdidos, vazamentos de
  memória, estouros de buffer e falhas de segmentação são responsáveis por uma
  fração enorme dos bugs em programas C/C++.
\end{itemize}

\textbf{Princípios de uso seguro:}
\begin{enumerate}[noitemsep]
  \item Inicialize sempre (com \texttt{NULL} ou endereço válido).
  \item Use \texttt{const} para indicar imutabilidade (princípio do menor privilégio).
  \item Verifique limites antes de qualquer aritmética de ponteiros.
  \item Para \texttt{void *}, sempre faça cast antes de desreferenciar.
\end{enumerate}

\textbf{Próximos passos:} no Capítulo~8 estudaremos strings em profundidade --
uma aplicação natural de \texttt{char *}. No Capítulo~10 veremos estruturas
(\texttt{struct}), e no Capítulo~12, alocação dinâmica e estruturas de dados
encadeadas. Tudo se apoia nos ponteiros que você acabou de aprender.

\end{tcolorbox}

\begin{boaprática}
  Linguagens modernas (Java, C\#, Python, JavaScript) escondem ponteiros sob
  abstrações como \textit{referências} e \textit{coleta de lixo automática}.
  Mas \textbf{ponteiros estão lá}, escondidos. Quem entende ponteiros em C
  entende como essas linguagens funcionam internamente. Esse é um dos maiores
  benefícios de aprender C antes (ou junto) de outras linguagens.
\end{boaprática}

\end{document}
